\chapter{data-driven wind estimation and compensation}

This chapter describes how the multi-layer perceptron (MLP) and the long short-term memory (LSTM) network introduced in Section~\ref{sec:2.6ArtificialNN} are implemented and trained for the wind estimation task considered in this thesis. 
Section~\ref{sec:problem_formulation} defines the learning task, Sections~\ref{sec:data_basis} and~\ref{sec:data_preprocessing} describe the data and its preparation, Sections~\ref{sec:architectures} to~\ref{sec:hyperparameter_opt} cover the model architectures, the training procedure and the hyperparameter search, and Section~\ref{sec:evaluation_metrics} introduces the metrics used for the comparison in Chapter~\ref{chap:results}.



\section{Problem Formulation}
\label{sec:4.1ProblemFormulation}

The task of the networks is to estimate the wind acting on the airship at a given time step $t$ from the vehicle's state and sensor information at or up to this time step.

To achieve this, the networks are provided with the same information available to the Kalman Filter from the onboard sensor suite described in Section~\ref{sec:1.4.1Hardware}, with one intentional exception. The position is left out, as the wind does not depend on the location of the airship. The position measurement therefore carries no information about the wind and could only serve to distinguish individual training runs, which promotes overfitting. The input vector is given by

\begin{equation}
	\bm{x}_t = \begin{bmatrix} \bm{v}_t^T & \bm{\Phi}_t^T & \bm{\omega}_t^T & \bm{a}_t^T & \bm{\dot{\omega}}_t^T & \bm{u}_t^T \end{bmatrix}^T \in \mathbb{R}^{24},
	\label{eq:nn_input}
\end{equation}

comprising velocity $\bm{v}$, attitude $\bm{\Phi}$, and angular rate $\bm{\omega}$, the measured linear accelerations $\bm{a}$ and angular accelerations $\bm{\dot{\omega}}$, and the motor inputs $\bm{u} \in \mathbb{R}^{9}$.
The velocity, angular rate, and accelerations are expressed in body-fixed coordinates. This yields an input vector dimension of $n = 24$.

The airspeed measurement of the pitot tube is not included in the input vector for simplicity. The sensor only delivers valid readings above a cutoff airspeed and outputs zero below it, so the measurement is not continuously available. Feeding such an intermittent signal to the networks would require additional handling, e.g., an availability flag, which is beyond the scope of this work. 

The network output is the wind velocity vector, $\hat{\bm{y}}_t = {}_E \hat{\bm{v}}_{W, t} \in \mathbb{R}^{3}$. In contrast to the KF, which estimates the wind in Earth-fixed coordinates, the networks estimate the wind in body-fixed coordinates. This choice simplifies the learning task for two reasons. First, the wind affects the airship dynamics through the aerodynamic velocity, which is defined in the body frame. Hence, $\bm{v}^{B}_{w}$ is the quantity that directly enters the aerodynamic forces and moments. Second, the networks do not have to learn the coordinate transformation between the body and the Earth frame.

With these definitions, both networks realize the mapping \eqref{eq:nn_mapping} with $n = 24$ and $m = 3$. The MLP maps the input $\bm{x}_t$ of a single time step to $\hat{\bm{y}}_t$, whereas the LSTM additionally processes a sequence of preceding inputs, as described in Section~\ref{sec:4.2.3InputGeneration}.





\section{Data Set}
\label{sec:4.2DataSet}

Due to the lack of real flight data, the networks are trained and evaluated using simulated data. The data set consists of 486 simulation runs, obtained by combining six flight scenarios (Curve, Forward fast, Forward slow, Hover, Rotate, and Sideways) with 81 wind scenarios each. The flight scenarios cover characteristic operating conditions of the airship, ranging from stationary hovering to translational and rotational maneuvers at different speeds.
Each wind scenario combines one of 27 distinct constant wind vectors with one of three turbulence levels, defined by $W_{20} \in \{0, 5, 10\}$\SI{}{m/s}, where $W_{20} = 0$ \SI{}{m/s} corresponds to the absence of turbulence.

Each simulation run lasts \SI{100}{\second}, with signals recorded at a sampling frequency of \SI{100}{\hertz}. In total, the data set covers \SI{48600}{\second}, or about 13.5 hours, of simulated flight. For each run, the recorded data consist of the input vector $\bm{x}_t$ defined in \eqref{eq:nn_input} and the corresponding simulated wind velocity in body-fixed coordinates, which serves as the target vector $\bm{y}_t$.





\subsection{Data Splitting}
\label{sec:4.2.1DataSplitting}

To evaluate the networks' generalization performance, the dataset is divided into separate training, validation, and test sets. The training set is used to update the network parameters, while the validation set is used to monitor training progress. The test set is reserved exclusively for the final evaluation. To prevent data leakage, each simulation run is assigned entirely to a single set.

Two different data-splitting strategies are considered to assess generalization to unseen simulation runs (Experiment A) and unseen flight scenarios (Experiment B).

For Experiment A, the data set is split into 342 training, 72 validation, and 72 test runs. The simulation runs are randomly assigned to the three sets separately for each combination of flight scenario and wind turbulence level. This ensures that all flight scenarios and turbulence levels are represented in each set.

For Experiment B, generalization to unseen flight scenarios is evaluated using six leave-one-scenario-out splits (B1--B6), each holding out one flight scenario entirely. All 81 runs of the held-out scenario form the test set. The runs from the remaining five scenarios are divided into 325 training and 80 validation runs, with the runs distributed as evenly as possible across the three turbulence levels.
Since the same wind scenarios are used across all flight scenarios, the test scenario differs from the training scenarios only in the airship's motion, while the wind conditions remain identical.

Table~\ref{tab:data_splits} provides an overview of the data splits used in the two experiments.

\begin{table}[htb]
	\centering
	\caption{Number of simulation runs in each data split. In Experiment B, the held-out flight scenario forms the test set.}
	\label{tab:data_splits}
	\begin{tabular}{lcccc}
		\toprule
		Experiment & Training & Validation & Test & Test data \\
		\midrule
		A (unseen runs)        & 342 & 72 & 72 & all scenarios \\
		B1--B6 (unseen scenario) & 325 & 80 & 81 & held-out scenario \\
		\bottomrule
	\end{tabular}
\end{table}





\subsection{Preprocessing}
\label{sec:4.2.2Preprocessing}

The sensor measurements are originally recorded at \SI{100}{\hertz}. Consecutive samples at this rate are highly redundant, and the resulting amount of data increases the computational cost considerably without adding much information. Therefore, all signals are downsampled to \SI{20}{\hertz} by retaining every fifth sample.

Furthermore, the inputs and targets are standardized to zero mean and unit variance, which improves the conditioning of the optimization problem and can accelerate training \cite{Lecun2000EffBackProp}. Each of the 24 components $x_{j,t}$ of the input vector $\bm{x}_t$ (e.g., a single velocity component or a single motor input) and each of the three components $y_{k,t}$ of the target vector $\bm{y}_t$ is standardized individually:

\begin{equation}
	\tilde{x}_{j,t} = \frac{x_{j,t} - \mu_{x,j}}{\sigma_{x,j}}, \quad j = 1, \dots, 24, \qquad
	\tilde{y}_{k,t} = \frac{y_{k,t} - \mu_{y,k}}{\sigma_{y,k}}, \quad k = 1, 2, 3,
	\label{eq:standardization}
\end{equation}

where $\mu$ and $\sigma$ denote the corresponding means and standard deviations. For each of the seven data splits, these statistics are computed using all time steps of all runs in the respective training set. They are then applied unchanged to the corresponding validation and test sets, preventing information from these sets from leaking into the training process.

The standardized input and target vectors are obtained by stacking their respective components:

\begin{equation}
	\tilde{\bm{x}}_t = [\tilde{x}_{1,t}, \dots, \tilde{x}_{24,t}]^T 
	\qquad
	\tilde{\bm{y}}_t = [\tilde{y}_{1,t}, \tilde{y}_{2,t}, \tilde{y}_{3,t}]^T.
\end{equation}

Both networks are trained using these standardized vectors. Standardizing the targets additionally prevents wind components with larger numerical magnitudes from dominating the loss function.
For readability, the tilde is omitted in the following, and $\bm{x}_t$ and $\bm{y}_t$ denote the standardized quantities unless stated otherwise.





\subsection{Input Generation}
\label{sec:4.2.3InputGeneration}

Based on the standardized and downsampled signals, the input-output pairs for the networks are assembled separately for each run. For the MLP, the input is the input vector $\bm{x}_t$ at the current time step, so no temporal information is provided.

For the LSTM, the input consists of a sequence of the most recent $T = 40$ input vectors,

\begin{equation}
	\bm{X}_t = \begin{bmatrix} \bm{x}_{t-T+1} & \dots & \bm{x}_{t-1} & \bm{x}_t \end{bmatrix}^T \in \mathbb{R}^{T \times 24},
	\label{eq:lstm_sequence}
\end{equation}

which corresponds to a time window of \SI{2}{\second} at a sampling frequency of \SI{20}{\hertz}. The sequences are generated using a sliding window with a stride of one time step. Both networks predict the wind $\hat{\bm{y}}_t$ at the current time step $t$. Each simulation run is treated as an independent time series, so the windows never cross run boundaries, and as the split is performed on the level of entire runs, no window overlaps two sets. For the LSTM, the first $T-1$ time steps of each run, for which no complete history is available, are not used as prediction targets. A run with $N$ time steps after downsampling therefore yields $N - T + 1$ sequences, whereas the MLP uses all $N$ time steps of each run.






\section{Model Architectures}
\label{sec:4.3ModelArchitectures}

This section describes the architectures of the two networks. Both map the input vectors defined in Section~\ref{sec:4.2.3InputGeneration} to an estimate of the body-fixed wind velocity and use a linear output layer, as is common for regression tasks. They differ in how the input is processed: the MLP evaluates the input of a single time step, whereas the LSTM processes a sequence of $T = 40$ input vectors. 

The architectures were selected manually rather than through a hyperparameter search. Since the networks are intended for real-time execution on onboard flight hardware, where computational resources and memory are limited, both are designed as moderately sized baseline models. Both achieve a low loss on the validation data (see Chapter~\ref{chap:5.0Evaluation}). However, since the validation data contain flight scenarios also represented in the training data, hyperparameter tuning based on the validation loss would primarily optimize performance on familiar scenarios. Its effect on generalization to unseen scenarios, as investigated in Experiment B, is uncertain. The results should therefore be interpreted as the performance of manually selected baseline architectures rather than as the best performance achievable with these network types.



\subsubsection{Multi-Layer Perceptron}

The MLP is a fully connected feedforward network as introduced in Section~\ref{sec:mlp}. It maps the input vector $\bm{x}_t \in \mathbb{R}^{24}$ at a single time step to the wind estimate $\hat{\bm{y}}_t \in \mathbb{R}^{3}$ through three hidden layers with 64, 64, and 32 neurons, i.e., $L = 4$ layers in the notation of \eqref{eq:mlp_forward}. The layer dimensions are therefore

\begin{equation}
	24 \rightarrow 64 \rightarrow 64 \rightarrow 32 \rightarrow 3 .
\end{equation}

All hidden layers use the ReLU activation function $\varphi(z) = \max(0, z)$, applied element-wise. The output layer is linear, i.e., $\varphi^{[4]}(\bm{z}) = \bm{z}$, so that the output can take any real value, as required for the regression of the standardized wind components.



\subsubsection{Long Short-Term Memory Network}

The LSTM network consists of two stacked LSTM layers with 64 hidden units each, followed by a fully connected output layer. Its input is the sequence $\bm{X}_t$ defined in \eqref{eq:lstm_sequence}. Let $\bm{h}^{[l]}_\tau$ denote the hidden state of LSTM layer $l \in \{1, 2\}$ after processing the $\tau$-th element of the input sequence, $\tau = 1, \dots, T$, computed according to \eqref{eq:lstm_forget}--\eqref{eq:lstm_hidden}. The first LSTM layer processes the input sequence and returns its hidden state at every time step. These hidden states form the input sequence of the second layer, which returns only its final hidden state at the last time step, $\bm{h}^{[2]}_T$, to the output layer,

\begin{equation}
	\hat{\bm{y}}_t = \bm{W}_\mathrm{out} \, \bm{h}^{[2]}_T + \bm{b}_\mathrm{out},
	\label{eq:lstm_output_layer}
\end{equation}

where $\bm{W}_\mathrm{out} \in \mathbb{R}^{3 \times 64}$ and $\bm{b}_\mathrm{out} \in \mathbb{R}^{3}$. The final hidden state corresponds to the last input vector of the sequence, $\bm{x}_t$, and summarizes the entire window. The network therefore realizes a sequence-to-one mapping, in which a window of $T = 40$ input vectors yields a single wind estimate for the current time step. The hidden and cell states are initialized with zeros at the beginning of every sequence, so that all windows are processed independently of each other. As for the MLP, the output layer is linear.

The number of trainable parameters follows from the layer dimensions. A fully connected layer with $n_\mathrm{in}$ inputs and $n_\mathrm{out}$ outputs has $n_\mathrm{out}(n_\mathrm{in} + 1)$ parameters, including biases. An LSTM layer with $n_\mathrm{h}$ hidden units has $4 n_\mathrm{h} (n_\mathrm{in} + n_\mathrm{h} + 1)$ parameters, since each of the three gates and the candidate cell state in \eqref{eq:lstm_forget}--\eqref{eq:lstm_candidate} has a weight matrix $\bm{W} \in \mathbb{R}^{n_\mathrm{h} \times n_\mathrm{in}}$, a recurrent weight matrix $\bm{U} \in \mathbb{R}^{n_\mathrm{h} \times n_\mathrm{h}}$, and a bias vector $\bm{b} \in \mathbb{R}^{n_\mathrm{h}}$. Table~\ref{tab:architectures} summarizes the architectures. 

\begin{table}[htb]
	\centering
	\caption{Layer structure and number of trainable parameters of the two networks.}
	\label{tab:architectures}
	\begin{tabular}{llrr}
		\toprule
		Model & Layer & Output size & Parameters \\
		\midrule
		MLP  & Dense, ReLU    & $64$          & 1\,600 \\
		& Dense, ReLU    & $64$          & 4\,160 \\
		& Dense, ReLU    & $32$          & 2\,080 \\
		& Dense, linear  & $3$           & 99 \\
		\cmidrule(l){2-4}
		& Total          &               & 7\,939 \\
		\midrule
		LSTM & LSTM (sequence output) & $64 \times T$ & 22\,784 \\
		& LSTM (last output)     & $64$          & 33\,024 \\
		& Dense, linear          & $3$           & 195 \\
		\cmidrule(l){2-4}
		& Total                  &               & 56\,003 \\
		\bottomrule
	\end{tabular}
\end{table}

The LSTM has considerably more trainable parameters than the MLP. Additionally, it evaluates $T = 40$ recurrent steps for each prediction, whereas the MLP requires a single forward pass. These differences in model size and computational requirements have to be taken into account when comparing their predictive performance in Chapter~\ref{chap:5.0Evaluation}.





\section{Training Procedure}
\label{sec:4.4TrainingProcedure}

Both networks are implemented and trained using the Deep Learning Toolbox of MATLAB 2025b. They are trained as regression models on the standardized input and target vectors from Section~\ref{sec:4.2.2Preprocessing}, using the training set of the respective data split. Training is performed separately for each of the seven splits and both network types, which results in 14 trained models.

The networks are trained by minimizing the regression loss between the predicted and true standardized wind vectors. MATLAB uses a half-mean-squared-error loss for the MLP. For the LSTM's sequence-to-one regression task, the loss is one-half of the sum of squared errors across the three output components, without normalization by the number of responses. Both loss definitions are proportional to the conventional mean squared error in Eq.~\eqref{eq:mse_loss} and therefore have the same minimizer. The gradients are computed by backpropagation and, for the LSTM, backpropagation through time over the $T=40$ time steps of each sequence.

The parameters are updated with the Adam optimizer \cite{Kingma2014Adam} using the default decay rates of the implementation. The learning rate starts at $\eta_0 = 10^{-3}$ and is halved every 20 epochs. The training data is divided into mini-batches and shuffled before every epoch. The mini-batch size is 1024 for the MLP and 128 for the LSTM to account for the greater computational and memory requirements of processing sequences of 40 time steps.
Besides the L2 regularization of the implementation with its default factor of $10^{-4}$, no further regularization such as dropout is used. The weights are initialized with the default initializers of the implementation, i.e., the Glorot initialization for the input weights, orthogonal initialization for the recurrent weights of the LSTM, and zeros for the biases, except for the forget gate bias, which is initialized with one.

The loss on the validation set is evaluated every 50 iterations. Training stops after a maximum of 100 epochs, or earlier if the validation loss has not fallen below its smallest value so far in 30 consecutive evaluations. The network of the final iteration is used for the evaluation. Table~\ref{tab:training_settings} summarizes the settings.

\begin{table}[htb]
	\centering
	\caption{Training settings of the two networks.}
	\label{tab:training_settings}
	\begin{tabular}{lcc}
		\toprule
		& MLP & LSTM \\
		\midrule
		Loss                        & Half mean squared error & Half sum of squared errors \\
		Optimizer                   & \multicolumn{2}{c}{Adam} \\
		Initial learning rate       & \multicolumn{2}{c}{$10^{-3}$} \\
		Learning rate schedule      & \multicolumn{2}{c}{Halved every 20 epochs} \\
		L2 regularization           & \multicolumn{2}{c}{$10^{-4}$} \\
		Maximum number of epochs    & \multicolumn{2}{c}{100} \\
		Validation frequency        & \multicolumn{2}{c}{50 iterations} \\
		Early stopping patience     & \multicolumn{2}{c}{30 validations} \\
		Mini-batch size             & 1024 & 128 \\
		\bottomrule
	\end{tabular}
\end{table}

The performance of the trained networks is evaluated in Chapter~\ref{chap:5.0Evaluation}. As the networks are trained on standardized targets, their outputs are transformed back to physical units by inverting \eqref{eq:standardization} to obtain the wind estimates. 

